\documentclass[11pt]{article}
\usepackage[utf8]{inputenc}
\usepackage[T1]{fontenc}
\usepackage{amsmath,amssymb,amsthm}
\usepackage{booktabs}
\usepackage[margin=1in]{geometry}
\usepackage{hyperref}

\newtheorem{theorem}{Theorem}
\newtheorem{remark}{Remark}

\title{The number of two-sided classes of the $0$-Hecke monoid is not $p(n)$:\\
a refutation of conjecture \texttt{00000003843}}
\author{Submission to The Last Math Competition}
\date{}

\begin{document}
\maketitle

\begin{abstract}
Conjecture \texttt{00000003843} asserts that the number of two-sided classes of
the $0$-Hecke monoid $H_n(0)$ equals the partition number $p(n)$, and that the
number of left classes inside each two-sided class equals the number of distinct
hook lengths of the corresponding partition.  We show that $H_n(0)$ is
$\mathcal{J}$-trivial, so it has exactly $n!$ two-sided classes and $n!$ left
classes.  At $n = 3$ this gives $6$ two-sided classes against $p(3) = 3$, and
$6$ left classes against $\sum_{\lambda \vdash 3}
\#\{\text{distinct hook lengths of }\lambda\} = 8$.  Both halves of the
conjecture therefore fail.  The second half already fails at $n = 2$ and is
false under the alternative (Kazhdan--Lusztig) reading of the statement as well.
Everything is a finite computation; it is formalised in Lean~4 using only the
kernel and core, and every computational theorem is verified to depend on no
axioms.
\end{abstract}

\section{The conjecture}

The conjecture reads, verbatim:

\begin{quote}
\textbf{Definition}: The $0$-Hecke monoid $H_n(0)$ is generated by monoid braid
relations and absorption relations.  \textbf{Conjecture}: The number of two-sided
classes of $H_n(0)$ is exactly the partition number $p(n)$, and the number of
left classes in each two-sided class equals the number of distinct hook lengths
of the corresponding partition. \emph{(0-Hecke two-sided class-partition)}
\end{quote}

Throughout, $H_n(0)$ is the monoid generated by $\pi_1,\dots,\pi_{n-1}$ subject
to the braid relations, far-commutation, and the absorption relations
$\pi_i^2 = \pi_i$.  Its elements are indexed by permutations $w \in S_n$, so
$\lvert H_n(0)\rvert = n!$, and multiplication by a generator is given by the
right-descent rule
\[
  \pi_w \cdot \pi_i =
  \begin{cases}
    \pi_{w s_i}, & \ell(w s_i) > \ell(w),\\
    \pi_w,       & \text{otherwise,}
  \end{cases}
\]
where $w s_i$ is $w$ with the entries in positions $i$ and $i+1$ swapped.
``Two-sided classes'' and ``left classes'' are read as Green's $\mathcal{J}$- and
$\mathcal{L}$-classes: $a \mathrel{\mathcal{J}} b$ iff $HaH = HbH$, and
$a \mathrel{\mathcal{L}} b$ iff $Ha = Hb$.  For a finite monoid
$\mathcal{D} = \mathcal{J}$, so this is also the two-sided/left cell reading
standard in semigroup theory.  Section~\ref{sec:kl} treats the alternative
Kazhdan--Lusztig reading.

\section{The computation}

Building $H_n(0)$ for $n \le 5$ and computing the principal ideals by closing
singletons under left and right multiplication by the generators gives
Table~\ref{tab:green}.  Three self-checks are included: the multiplication is
verified to be associative on all triples (an exhaustive check, not a sample),
the identity element acts as a two-sided identity, and the number of idempotents
comes out as $2,4,8,16 = 2^{n-1}$, which is the known count for $H_n(0)$.

\begin{table}[htbp]
\centering
\caption{Green's classes of $H_n(0)$ against the conjectured values.}
\label{tab:green}
\begin{tabular}{@{}rrrrrr@{}}
\toprule
$n$ & $\lvert H_n(0)\rvert$ & $\mathcal{J}$-classes & $\mathcal{L}$-classes
    & $p(n)$ & $\displaystyle\sum_{\lambda\vdash n}
               \#\{\text{distinct hook lengths}\}$\\
\midrule
2 &   2 &   2 &   2 & 2 &  4\\
3 &   6 &   6 &   6 & 3 &  8\\
4 &  24 &  24 &  24 & 5 & 17\\
5 & 120 & 120 & 120 & 7 & 29\\
\bottomrule
\end{tabular}
\end{table}

The number of $\mathcal{J}$-classes equals the number of elements, i.e.\ every
principal two-sided ideal is distinct and $H_n(0)$ is $\mathcal{J}$-trivial (and
hence also $\mathcal{L}$- and $\mathcal{R}$-trivial).  This is consistent with
the literature: Ganyushkin and Mazorchuk prove that $H_\Gamma$ is
$\mathcal{J}$-trivial \cite{ganyushkin2010}.

For $n = 3$ the six principal two-sided ideals are, as characteristic vectors on
$(\mathrm{id}, s_1, s_2, s_1s_2, s_2s_1, w_0)$:
\[
\begin{aligned}
J(\mathrm{id})   &= (1,1,1,1,1,1), & J(s_1)   &= (0,1,0,1,1,1),\\
J(s_2)           &= (0,0,1,1,1,1), & J(s_1s_2)&= (0,0,0,1,0,1),\\
J(s_2s_1)        &= (0,0,0,0,1,1), & J(w_0)   &= (0,0,0,0,0,1).
\end{aligned}
\]
They are pairwise distinct, so there are $6$ two-sided classes.  Note the clean
structural reason: $\mathcal{J}$-triviality means each $\mathcal{J}$-class is a
singleton, so each two-sided class contains exactly \emph{one} left class.

\begin{theorem}\label{thm:one}
The first half of conjecture \texttt{00000003843} is false: $H_3(0)$ has $6$
two-sided classes but $p(3) = 3$.
\end{theorem}

\begin{theorem}\label{thm:two}
The second half of conjecture \texttt{00000003843} is false: $H_3(0)$ has $6$
left classes in total, whereas the conjecture predicts
$\sum_{\lambda\vdash 3}\#\{\text{distinct hook lengths of }\lambda\}
= 3 + 2 + 3 = 8$.
\end{theorem}

\begin{remark}
The two counts in Theorem~\ref{thm:one} happen to agree at $n = 2$
($2 = p(2)$), so $n = 3$ is the smallest witness for the first half.  The second
half already fails at $n = 2$: there each $\mathcal{J}$-class is a singleton and
so contains one left class, while both partitions of $2$ have hook lengths
$\{2,1\}$ and hence $2$ distinct hook lengths.  Thus $1 \ne 2$ at $n = 2$.
\end{remark}

\section{The alternative (Kazhdan--Lusztig) reading}
\label{sec:kl}

The phrase ``the corresponding partition'' suggests that the generator had the
Kazhdan--Lusztig cell decomposition in mind: for the Hecke algebra of $S_n$ the
two-sided cells \emph{are} indexed by partitions, so their number is $p(n)$, and
the left cells inside the two-sided cell $C_\lambda$ are indexed by the standard
Young tableaux of shape $\lambda$, so there are $f^\lambda$ of them.

Under that reading the first half of the conjecture is true (for $S_n$, not for
the $0$-Hecke monoid) but the second half is still false, because
$f^\lambda$ is the number of standard Young tableaux and not the number of
distinct hook lengths.  At $n = 2$ we have $f^{(2)} = f^{(1,1)} = 1$, whereas
both partitions have hook lengths $\{2,1\}$, i.e.\ $2$ distinct hook lengths.
Already at $n = 2$ the predicted number is $2$ while the true number is $1$.

So the second half of the conjecture fails under \emph{both} readings, with the
same small witness, and the first half fails under the reading forced by the
subject of the statement (a monoid, hence Green's relations).

\section{Formalisation}

The directory \texttt{lean4/} contains a Lean~4 project using only the kernel
and core (no Mathlib, no \texttt{sorry}).  The file \texttt{Main.lean} defines
\begin{itemize}
  \item $S_3$ in one-line notation together with the inversion-number length,
        the generators $s_1, s_2$, and the right-descent product of $H_3(0)$;
  \item the principal two-sided and left ideals as closures under multiplication
        by generators, encoded as characteristic vectors;
  \item integer partitions, hook lengths, and the count of distinct hook lengths.
\end{itemize}
The theorems proved are
\[
\begin{array}{ll}
\texttt{hSize\_eq}              & \lvert H_3(0)\rvert = 6,\\
\texttt{assoc\_ok}              & \text{associativity on all }6^3\text{ triples},\\
\texttt{identity\_ok}           & \text{two-sided identity},\\
\texttt{absorption\_ok}         & \pi_i^2 = \pi_i\text{ for both generators},\\
\texttt{jClasses\_eq}           & \#\mathcal{J}\text{-classes} = 6,\\
\texttt{lClasses\_eq}           & \#\mathcal{L}\text{-classes} = 6,\\
\texttt{pCount3\_eq}            & p(3) = 3,\\
\texttt{sumHooks3\_eq}          & \sum_{\lambda\vdash 3}\#\{\text{distinct hook lengths}\} = 8,\\
\texttt{refutation\_one\_bool}  & (\#\mathcal{J}\text{-classes} \stackrel{?}{=} p(3)) = \mathrm{false},\\
\texttt{refutation\_two\_bool}  & (\#\mathcal{L}\text{-classes} \stackrel{?}{=} \text{hook sum}) = \mathrm{false},\\
\texttt{refutation\_part\_one}  & \text{Theorem~\ref{thm:one}},\\
\texttt{refutation\_part\_two}  & \text{Theorem~\ref{thm:two}}.
\end{array}
\]
Running \texttt{lake env lean Check.lean} prints every computed value and
\texttt{\#print axioms} for every theorem.  \textbf{All eight computational
theorems and both Boolean witnesses depend on no axioms at all.}  The two
propositional statements \texttt{refutation\_part\_one} and
\texttt{refutation\_part\_two} (and the conjunction \texttt{summary}) depend on
\texttt{[propext]} only, which enters through Lean's generated no-confusion
machinery for $\lnot$; no \texttt{Quot.sound}, \texttt{Classical.choice},
\texttt{sorryAx} or Mathlib axiom is used anywhere.

One implementation note is worth recording, since it is what makes the
axiom-free result possible: Lean~4's short-circuiting \texttt{\&\&} and
\texttt{||} elaborate to \texttt{if} on a \texttt{Bool}, which goes through
\texttt{CoeSort Bool Prop} and \texttt{Decidable (b = true)} and hence pulls in
\texttt{propext}; and the standard-library list accessor with a default value
does the same.  Both are therefore replaced here by functions defined directly
by pattern matching (\texttt{band}, \texttt{bor}, \texttt{idxD}).

\section{Reproduction}

\texttt{python3 reproduce.py} (standard library only, no third-party packages)
rebuilds $H_n(0)$ for $n = 2,\dots,5$, re-runs the three self-checks, recomputes
Table~\ref{tab:green}, prints the six principal two-sided ideals of $H_3(0)$ and
the hook-length bookkeeping for the partitions of $3$.  Its output agrees
cell-by-cell with the Lean computation.

\begin{thebibliography}{9}
\bibitem{ganyushkin2010}
O.~Ganyushkin and V.~Mazorchuk,
\emph{On Kiselman quotients of $0$-Hecke monoids},
Int.\ Electron.\ J.\ Algebra \textbf{10} (2011), 174--204;
arXiv:1006.0316.  (The $\mathcal{J}$-triviality of $H_\Gamma$ is derived there.)
\end{thebibliography}

\end{document}
